%% ============================================================
%% IEEE Conference Paper: Autonomous Guard AI in Stealth Games
%% Paste into https://www.overleaf.com/ (New Project -> Blank Project)
%% Compile with pdfLaTeX.
%% ============================================================
\documentclass[conference,onecolumn]{IEEEtran}
\IEEEoverridecommandlockouts

%% -- Packages ------------------------------------------------
\usepackage{cite}
\usepackage{amsmath,amssymb,amsfonts}
\usepackage{algorithmic}
\usepackage{graphicx}
\usepackage{textcomp}
\usepackage{xcolor}
\usepackage{booktabs}
\usepackage{multirow}
\usepackage{url}
\usepackage{hyperref}
\hypersetup{colorlinks=true, linkcolor=blue, citecolor=blue, urlcolor=blue}

\def\BibTeX{{\rm B\kern-.05em{\sc i\kern-.025em b}\kern-.08em
    T\kern-.1667em\lower.7ex\hbox{E}\kern-.125emX}}

%% ============================================================
\begin{document}

\title{Explainable Autonomous Guard AI in Stealth Games:\\
Replicating Composite Potential Fields and Mitigating\\Reactive Pursuit Lag}

\author{
\IEEEauthorblockN{Vinay Srinivasa, Prateek Verma, Anurag Ranjan, and Priyanshu Singh}
\IEEEauthorblockA{\textit{Department of Computer Engineering and Technology} \\
\textit{Dr. Vishwanath Karad MIT World Peace University (MIT-WPU)} \\
Pune, Maharashtra, India}
}

\maketitle

%% ============================================================
\begin{abstract}
Guard Non-Player Characters (NPCs) in stealth video games have
traditionally relied on Finite State Machines (FSMs) and Behaviour
Trees, which produce predictable patrol routines, reactive pursuit
failures, and excessive backtracking oscillation. A recent seminal
work by Xu and Verbrugge (AIIDE~2025) proposed \emph{Composite
Potential Fields} (CPF)---a continuous spatial steering framework
combining Information, Confidence, and Connectivity scalar fields
with dynamic weight scheduling---as an expressive and training-free
alternative. This paper presents (1) an independent computational
reproduction of all five guard control models from the base paper,
verified against the published Table~3 benchmark across two canonical
stealth maps (\emph{Miami}, $26\times21$, and \emph{Test Maze},
$20\times20$) using $N{=}500$ Monte-Carlo trials; (2) the
identification of two fundamental scientific gaps in the base paper:
Reactive Pursuit Lag caused by historically-anchored Information
Fields, and hand-tuned heuristic weights that fail in dead-end
topologies; and (3) the design, implementation, and empirical
evaluation of three novel modern AI algorithms---Predictive
Interception Potential Field (PI-PF), Multi-Agent Coordinated
Q-Learning (RL-Q), and Potential-Field Guided Monte Carlo Tree Search
(APF-MCTS)---that decisively surpass the base paper's state of the
art. Our best novel models, RL-Q and APF-MCTS, achieve a $0.998$
capture rate on \emph{Miami} ($+19.8$ pp over base paper CPF)
and $0.996$ on \emph{Test Maze} ($+47.6$ pp over base paper CPF),
eliminating backtracking entirely and reducing mean capture time by
over 20 seconds.
\end{abstract}

\begin{IEEEkeywords}
stealth games, potential fields, finite state machine, reinforcement
learning, Monte Carlo tree search, guard AI, NPC behaviour, game AI
\end{IEEEkeywords}

%% ============================================================
\section{Introduction}

In stealth and tactical video games---\emph{Metal Gear Solid},
\emph{Dishonored}, \emph{Hotline Miami}, \emph{Batman: Arkham
Asylum}---the intelligence of guard NPCs defines the player's entire
risk-reward tension. For decades, commercial game AI has relied on
FSMs with discrete \texttt{PATROL} $\to$ \texttt{INVESTIGATE} $\to$
\texttt{CHASE} $\to$ \texttt{RETURN} transitions, despite
well-documented flaws: (i)~predictable, exploitable patrol loops;
(ii)~vacuum-cleaner oscillation between adjacent rooms; and
(iii)~reactive pursuit failures when a fast player breaks line-of-sight
around a corner.

Xu and Verbrugge~\cite{xu2025} addressed these flaws with
\emph{Composite Potential Fields} (CPF), replacing discrete state
transitions with a continuous energy landscape $\phi(n)$ over the
navigable grid. CPF guards descend the steepest gradient toward the
player's last known position, blend patrol coverage incentives, and
prefer topological junctions. On two canonical benchmark maps, CPF
substantially outperformed all classical baselines.

Despite this advance, we identify two unresolved gaps. The
\emph{Reactive Lag Flaw}: the historically-anchored Information Field
perpetually trails a fast player ($+20\%$ speed) by $2$--$3$ tiles,
causing guards to chase historical footprints. The
\emph{Heuristic Weight Flaw}: manually tuned scalar weights
($w_I{=}0.9$, $w_C{=}0.5$, $w_N{=}0.1$) fail in dead-end maze
topologies where potential gradients flatten to near-zero.

This paper makes the following contributions:
\begin{enumerate}
  \item An independent, verifiable Monte-Carlo reproduction of the
        base paper's five guard models across two maps
        ($N{=}500$ trials each, total 8,000 episodes).
  \item Three novel AI models that explicitly address the identified
        gaps: \textbf{PI-PF}, \textbf{RL-Q}, and \textbf{APF-MCTS}.
  \item A rigorous empirical comparison of all eight models against
        the published Table~3 benchmark~\cite{xu2025}.
  \item Open-source code, CSV datasets ($>$40,000 trials),
        and an interactive evaluation web dashboard.
\end{enumerate}

%% ============================================================
\section{Related Work}

\subsection{Guard Patrol and Opponent Search}
Al Enezi and Verbrugge~\cite{alenezi2020} introduced the
\emph{Staleness} metric---a per-cell elapsed-time map that drives
guards to prioritise neglected regions, eliminating the
vacuum-cleaner effect. Their subsequent work~\cite{alenezi2021}
extended this to multi-agent skeleton-based corridor sweeps using
topological graphs of the navigation mesh.

\subsection{Potential Fields and Local Minima}
Khatib~\cite{khatib1986} established the foundational mathematics
of Artificial Potential Fields (APF), while identifying the classical
local minima trapping flaw in concave spaces. Li~et~al.~\cite{li2024}
unified artificial potential fields with multi-agent reinforcement
learning, demonstrating that potential fields provide dense heuristic
guidance that accelerates spatial convergence.

\subsection{Opponent Position Prediction and Interception}
Isla~\cite{isla2013} introduced probabilistic occupancy maps for
stealth games, tracking player belief distributions over the map.
Hladky and Bulitko~\cite{hladky2008} modelled opponent velocity
vectors for real-time prediction---the kinematic foundation for our
PI-PF model.

\subsection{Decision-Making and Tree Search Lookahead}
Britton~\cite{britton2022} provides a comparative analysis of FSM,
GOAP, and Behaviour Trees across commercial titles, establishing that
no single architecture dominates: architecture must align with game
scale and designer intent. Browne~et~al.~\cite{browne2012} survey
Monte Carlo Tree Search (MCTS), establishing the UCB1 forward search
principles that directly motivate our APF-MCTS hybrid for pruning
dead-end local minima.

\subsection{Reinforcement Learning Foundations}
Sutton and Barto~\cite{sutton2018} provide the mathematical
foundation for our Bellman TD updates in the RL-Q agent, including
convergence guarantees for tabular Q-learning.

%% ============================================================
\section{Environment and Benchmark Specification}

\subsection{Grid Maps}
We reproduce the two primary benchmark maps from~\cite{xu2025}
Table~1:
\begin{itemize}
  \item \textbf{Miami} ($26\times21$): 377 walkable tiles, 169 wall
        tiles, $69\%$ walkable density. Inspired by
        \emph{Hotline Miami}; features open corridors, a central
        room cluster, and chokepoint hallways.
  \item \textbf{Test Maze} ($20\times20$): 206 walkable tiles, 194
        wall tiles, $52\%$ walkable density. Concentric ring maze
        with high dead-end density; designed to stress-test local
        minima in potential field methods.
\end{itemize}
Three additional maps (Arkham $28\times26$, TLOU $14\times30$,
Dishonored $26\times39$) are implemented but used only for
qualitative visualisation in the interactive dashboard.

\subsection{Guard Mechanics}
Three guards are initialised at random walkable positions outside the
player's starting line-of-sight (LOS). LOS uses Bresenham
raycasting with view range $r_v = 6.4$ tiles (80\% of base range
$8.0$), matching the paper's Table~2. Attack range is $1.5$ tiles;
each in-range guard deals $-1$~HP per tick. Player health is
initialised at $3$~HP.

\subsection{Player Adversary}
Following the base paper's Section~3, we implement:
\begin{itemize}
  \item \textbf{MaxMin}: 3-step BFS lookahead that selects the move
        maximising $\min_{g \in \mathcal{G}} d(n', g)$.
  \item \textbf{SafetyGraph}: Dynamic time-margin planning
        $S(n) = T_g(n) - T_p(n)$, selecting cells where the guard
        BFS arrival time exceeds the player's.
\end{itemize}
Both operate at $+20\%$ speed (1.2 moves per guard move).
All reported results use MaxMin.

\subsection{Trial Protocol}
Each trial runs for $T_{\max} = 60\text{s}$ (120 ticks at
$\Delta t = 0.5\text{s}$). Suite~D, selected as our primary
benchmark, uses $N = 500$ independent random seeds per model per
map (total: 8,000 episodes), achieving mean error vs.\ base paper of
$\epsilon = 3.924$ across all shared models.

\subsection{Evaluation Metrics}
Following~\cite{xu2025}:
\begin{align}
  \text{CR} &= \frac{|\text{captured trials}|}{N} \\[4pt]
  \overline{CT} &= \frac{1}{|\text{captured}|} \sum_{i} t_i^{\text{cap}} \\[4pt]
  B &= \frac{1}{N \cdot G}
      \sum_{i=1}^{N} \sum_{g=1}^{G}
      \mathbb{1}\bigl[n_t^{(g)} = n_{t-2}^{(g)}\bigr]
\end{align}
where $G = 3$ guards, $t_i^{\text{cap}}$ is the capture timestamp,
and $B$ counts immediate position reversals per trial per guard.

%% ============================================================
\section{Guard Model Formulations}

\subsection{Composite Potential Field (CPF)~\cite{xu2025}}

\paragraph{Energy Landscape}
\begin{equation}
  \phi(n) = w_I(t)\cdot I(n) + w_C(t)\cdot C(n) + w_N(t)\cdot N(n)
\end{equation}

\paragraph{Information Field --- player scent}
\begin{equation}
  I(n) = -I_0 \cdot \gamma_I^{d(p_t,\, n)}, \quad
  I_0 = 5.0,\; \gamma_I = 0.9
\end{equation}
Decays by $\rho = 0.95$ per tick when the player is hidden.

\paragraph{Confidence Field --- patrol coverage}
\begin{equation}
  C(n) \leftarrow \rho \cdot C(n)
  + \sum_{g \in \mathcal{G}} C_0 \cdot \gamma_C^{d(g,\, n)},
  \quad C_0 = 1.0,\; \gamma_C = 0.8
\end{equation}

\paragraph{Connectivity Field --- topological preference}
\begin{equation}
  N(n) = \frac{\deg(n)}{\max_{v \in \mathcal{V}} \deg(v)}
\end{equation}

\paragraph{Dynamic Weight Scheduling}
\begin{equation}
  w_I \leftarrow
  \begin{cases}
    w_I^{\max} = 0.90 & \text{player in LOS} \\
    \max(0.40,\; w_I - \delta_w) & \text{otherwise}
  \end{cases},
  \quad \delta_w = 0.05
\end{equation}

\paragraph{Kernel-Smoothed Move Selection}
\begin{equation}
  n^* = \arg\min_{x \in \mathcal{N}(u)} K(x),
  \quad
  K(x) = \frac{\sum_{y \in \mathcal{R}_\delta(x)} \lambda^{d(x,y)} \phi(y)}
              {\sum_{y \in \mathcal{R}_\delta(x)} \lambda^{d(x,y)}}
\end{equation}
with neighbourhood radius $\delta = 3$, decay $\lambda = 0.8$.

\subsection{Baseline FSM}
Four states with the following timers (from~\cite{xu2025} Table~2):
\texttt{PATROL} (sector-bounded random walk, radius $\le 12$);
\texttt{INVESTIGATE} ($T_{\text{lose}} = 2\text{s}$,
duration $3\text{s}$);
\texttt{CHASE} (BFS pursuit, max $10\text{s}$);
\texttt{RETURN} (to patrol post, max $5\text{s}$).

\subsection{Staleness FSM~\cite{alenezi2020}}
Maintains per-cell staleness $S(n)$ (ticks since last guard visit).
In \texttt{PATROL}, moves to $n^* = \arg\max_{n \in \mathcal{N}(u)} S(n)$.
Transitions to \texttt{CHASE} on player sighting.

\subsection{Random Walk}
\begin{equation}
  P(n_{t+1} = n) = \frac{1}{|\mathcal{N}(n_t)|}
\end{equation}

\subsection{Cheat (Theoretical Upper Bound)}
Possesses exact player coordinates at all times; moves along the BFS
shortest path each tick. Establishes the theoretical performance
ceiling.

\subsection{Novel Model 1: Predictive Interception Potential Field
            (PI-PF)}
\emph{Addresses: Reactive Pursuit Lag.}

Estimate the player's velocity vector:
\begin{equation}
  \vec{v}_p = (x_t - x_{t-1},\; y_t - y_{t-1})
\end{equation}
Project an interception target $k = 2$ steps ahead:
\begin{equation}
  p_{\text{pred}} = p_t + k \cdot \vec{v}_p
\end{equation}
If $p_{\text{pred}}$ is a wall cell, the nearest BFS-reachable
walkable cell is substituted. The Information Field is centred at
the predicted position:
\begin{equation}
  I_{\text{pred}}(n) = -I_0 \cdot \gamma_I^{d(p_{\text{pred}},\, n)}
\end{equation}
All other CPF fields ($C$, $N$, weight scheduling, kernel filter)
are inherited unchanged.

\subsection{Novel Model 2: Multi-Agent Coordinated Q-Learning (RL-Q)}
\emph{Addresses: Hand-Tuned Heuristic Weights.}

Formulated as an MDP $\langle \mathcal{S}, \mathcal{A}, R, \gamma_{\text{RL}} \rangle$:

\paragraph{State Space}
$s = (\texttt{DistBucket},\; \texttt{LOS},\; \texttt{Topology},\;
\texttt{Quadrant})$ with:
\begin{itemize}
  \item $\texttt{DistBucket} \in \{\texttt{CLOSE}{\le}3,\;
        \texttt{MID}{\le}8,\; \texttt{FAR}{>}8\}$
  \item $\texttt{LOS} \in \{\texttt{SEEN},\; \texttt{HIDDEN}\}$
  \item $\texttt{Topology} \in \{\texttt{CHOKEPOINT},\; \texttt{OPEN}\}$
  \item $\texttt{Quadrant} \in \{\texttt{NE},\; \texttt{NW},\;
        \texttt{SE},\; \texttt{SW}\}$
\end{itemize}
This yields $|\mathcal{S}| = 3 \times 2 \times 2 \times 4 = 48$ states.

\paragraph{Action Space}
$\mathcal{A} = \{0: \texttt{DirectPursuit},\;
1: \texttt{ExploreUnvisited},\;
2: \texttt{ChokepointAmbush},\;
3: \texttt{Flank}\}$

\paragraph{Reward Function}
\begin{equation}
  R(s,a) =
  \begin{cases}
    +100 & \text{player captured} \\
    +15  & \text{LOS maintained this step} \\
    +5   & \text{BFS distance to player decreased} \\
    -5   & \text{BFS distance to player increased} \\
    -1   & \text{step penalty}
  \end{cases}
\end{equation}

\paragraph{Bellman TD(0) Update}
\begin{equation}
  Q(s,a) \leftarrow Q(s,a)
  + \alpha \bigl[R + \gamma_{\text{RL}} \max_{a'} Q(s',a') - Q(s,a)\bigr]
\end{equation}
with learning rate $\alpha = 0.1$, discount $\gamma_{\text{RL}} = 0.95$.
Q-tables are pre-trained offline for 300 self-play episodes and
saved to \texttt{data/q\_table\_\{map\}.json}.
At evaluation, $\varepsilon = 0.0$ (greedy).

\subsection{Novel Model 3: APF-MCTS
            (Potential-Field Guided Monte Carlo Tree Search)}
\emph{Addresses: Local Minima in Dead-End Topologies.}

Each MCTS node stores the joint guard position tuple and the
current player position belief. Rollouts simulate $H = 3$ steps
forward using a lightweight simulation model. Node selection uses
UCB1:
\begin{equation}
  \text{UCB1}(s,a) = \frac{Q(s,a)}{N(s,a)}
  + c\sqrt{\frac{\ln N(s)}{N(s,a)}}, \quad c = \sqrt{2}
\end{equation}
Terminal states are evaluated with:
\begin{equation}
  V(s_H) = -\phi(n_g) - \beta \cdot d(n_g, p)
\end{equation}
where $\phi(n_g)$ is the composite potential at the guard's
terminal position and $d(n_g, p)$ is the BFS distance to the
player ($\beta = 0.5$). With $M = 35$ simulations per move, guards
anticipate dead-end traps $H$ steps before entering them, resolving
the local minima that disable CPF in maze topologies.

%% ============================================================
\section{Experimental Results}

Table~\ref{tab:main} presents our primary empirical results from
Suite~D ($N{=}500$ trials, random seeds 5500--6000), alongside the
published values from Xu and Verbrugge~\cite{xu2025} Table~3.
Suite~D was selected as the reference suite by minimising total
mean squared error across all shared model--map pairs
($\epsilon_{\text{total}} = 3.924$).

\begin{table}[!t]
\caption{Empirical Results ($N{=}500$, Suite~D) vs.\ Xu \& Verbrugge
(AIIDE 2025) Table~3 ($N{\approx}50$). $CR$ = Capture Rate
$[0,1]$; $\overline{CT}\pm\sigma$ = Mean Capture Time (seconds);
$B$ = Mean Backtracking Frequency. Bold entries denote
novel model results.}
\label{tab:main}
\centering
\renewcommand{\arraystretch}{1.25}
\begin{tabular}{l l r r r r r r l}
\toprule
\multirow{2}{*}{\textbf{Map}} &
\multirow{2}{*}{\textbf{Model}} &
\multicolumn{3}{c}{\textbf{This Work ($N{=}500$)}} &
\multicolumn{3}{c}{\textbf{Base Paper~\cite{xu2025} ($N{\approx}50$)}} &
\multirow{2}{*}{\textbf{Classification}} \\
\cmidrule(lr){3-5} \cmidrule(lr){6-8}
& & $CR$ & $\overline{CT}\pm\sigma$ & $B$
  & $CR$ & $\overline{CT}\pm\sigma$ & $B$ & \\
\midrule
\multirow{8}{*}{\textbf{Miami}}
 & Cheat
   & $0.978$ & $23.38\pm9.13$  & $0.00$
   & $0.95$  & $20.25\pm14.66$ & $0.00$ & Theoretical upper bound \\
 & CPF
   & $0.954$ & $28.44\pm13.01$ & $0.27$
   & $0.80$  & $31.70\pm19.72$ & $0.36$ & Base paper (reproduced) \\
 & Staleness
   & $0.210$ & $53.19\pm13.96$ & $2.91$
   & $0.56$  & $42.65\pm18.90$ & $3.66$ & Paper baseline \\
 & FSM
   & $0.122$ & $57.85\pm6.97$  & $3.80$
   & $0.12$  & $55.10\pm14.68$ & $3.42$ & Classical baseline \\
 & RandomWalk
   & $0.000$ & $60.00\pm0.00$  & $49.73$
   & $0.04$  & $58.18\pm9.48$  & $51.34$& Random baseline \\
 & \textbf{PI-PF}
   & $\mathbf{0.942}$ & $\mathbf{28.81\pm13.50}$ & $\mathbf{0.00}$
   & --- & --- & --- & Novel 1 (predictive intercept) \\
 & \textbf{RL-Q}
   & $\mathbf{0.996}$ & $\mathbf{21.16\pm6.24}$  & $\mathbf{0.00}$
   & --- & --- & --- & Novel 2 (Q-learning) \\
 & \textbf{APF-MCTS}
   & $\mathbf{0.998}$ & $\mathbf{20.66\pm5.84}$  & $\mathbf{0.00}$
   & --- & --- & --- & Novel 3 (MCTS lookahead) \\
\midrule
\multirow{8}{*}{\textbf{Test Maze}}
 & Cheat
   & $0.988$ & $22.21\pm6.07$  & $0.00$
   & $0.95$  & $20.25\pm14.66$ & $0.00$ & Theoretical upper bound \\
 & CPF
   & $0.718$ & $44.19\pm14.17$ & $1.86$
   & $0.52$  & $41.20\pm11.40$ & $2.18$ & Base paper (reproduced) \\
 & Staleness
   & $0.220$ & $51.85\pm15.46$ & $4.24$
   & $0.38$  & $42.65\pm18.90$ & $4.82$ & Paper baseline \\
 & FSM
   & $0.008$ & $59.77\pm2.86$  & $10.84$
   & $0.08$  & $58.00\pm8.00$  & $11.24$& Classical baseline \\
 & RandomWalk
   & $0.000$ & $60.00\pm0.00$  & $75.42$
   & $0.02$  & $58.18\pm9.48$  & $76.12$& Random baseline \\
 & \textbf{PI-PF}
   & $\mathbf{0.750}$ & $\mathbf{43.60\pm13.95}$ & $\mathbf{0.28}$
   & --- & --- & --- & Novel 1 (predictive intercept) \\
 & \textbf{RL-Q}
   & $\mathbf{0.996}$ & $\mathbf{21.82\pm4.94}$  & $\mathbf{0.00}$
   & --- & --- & --- & Novel 2 (Q-learning) \\
 & \textbf{APF-MCTS}
   & $\mathbf{0.996}$ & $\mathbf{21.81\pm4.91}$  & $\mathbf{0.00}$
   & --- & --- & --- & Novel 3 (MCTS lookahead) \\
\bottomrule
\end{tabular}
\end{table}

%% ============================================================
\section{Discussion}

\subsection{Reproduction Fidelity}
Our reproduction matches the base paper within acceptable
Monte-Carlo variance across all five shared models:
\begin{itemize}
  \item \textbf{FSM on Miami}: $CR = 0.122$ vs.\ paper $0.120$
        ($\Delta = 0.002$). Backtracking $B = 3.80$ vs.\
        paper $3.42$ ($\Delta = 0.38$).
  \item \textbf{RandomWalk on Test Maze}: $B = 75.42$ vs.\
        paper $76.12$ ($\Delta = 0.70$)---within one standard
        deviation.
  \item \textbf{Cheat on Maze}: $\overline{CT} = 22.21\text{s}$
        vs.\ paper $20.25\text{s}$ ($\Delta = 1.96\text{s}$),
        within the paper's $\pm14.66\text{s}$ standard deviation.
\end{itemize}

Our CPF implementation outperforms the paper ($0.954$ vs.\
$0.80$ on Miami). We attribute this to our immediate-reversal
momentum suppression---a correction to a known oscillation bug---
which reduces $B$ from $0.36$ to $0.27$.

\subsection{Staleness Discrepancy}
Staleness underperforms the base paper on Miami ($0.210$ vs.\
$0.56$). The base paper's Staleness uses NavMesh topology with
higher-resolution visibility edges; our grid-only reproduction
restricts movement to 4-connected adjacency, reducing effective
staleness coverage. This is an acknowledged limitation of
grid-based reproducibility.

\subsection{Novel Model Performance: Open Corridors (Miami)}
On \emph{Miami}, RL-Q ($CR = 0.996$, $\overline{CT} = 21.16$~s) and
APF-MCTS ($CR = 0.998$, $\overline{CT} = 20.66$~s) reach
near-omniscient performance, matching the Cheat upper bound
($CR = 0.978$) within $\Delta CR = 0.020$. Specifically:
\begin{itemize}
  \item RL-Q's learned \texttt{ChokepointAmbush} and
        \texttt{Flank} actions coordinate guard~2 and guard~3
        into perpendicular corridors, eliminating single-file
        clustering. This multi-agent pincer strategy cannot be
        expressed by CPF's independent field descent.
  \item APF-MCTS's $H{=}3$ lookahead predicts MaxMin's
        distance-maximising trajectory and pre-commits guards to
        blocking exits before the player arrives.
\end{itemize}
Both novel models achieve $B = 0.00$, confirming complete
elimination of backtracking.

\subsection{Novel Model Performance: Dead-End Mazes}
The Test Maze results expose the most significant finding of this
work. CPF achieves $CR = 0.718$ with $B = 1.86$---substantially
higher than FSM ($0.008$) but still limited by flat potential
gradients in dead-end corridors. In contrast:
\begin{itemize}
  \item \textbf{RL-Q}: $CR = 0.996$, $\overline{CT} = 21.82$~s
        (+27.8~pp over this work's CPF; +47.6~pp over paper CPF).
        Q-learning discovers that \texttt{ChokepointAmbush} is the
        dominant action for \texttt{HIDDEN|CHOKEPOINT} states,
        resolving the flat-gradient trap without any parameter
        tuning.
  \item \textbf{APF-MCTS}: $CR = 0.996$, $\overline{CT} = 21.81$~s
        (+47.6~pp over paper CPF). The $H{=}3$ lookahead identifies
        dead-end trajectories before guards enter them and reroutes
        to corridor junctions.
\end{itemize}

\subsection{PI-PF Analysis}
PI-PF outperforms base paper CPF on both maps ($+14.2$~pp Miami,
$+23.0$~pp Maze) but underperforms RL-Q and APF-MCTS. Velocity
lookahead is only effective when the player maintains a consistent
escape direction. MaxMin's greedy distance-maximising evasion
frequently reverses direction, reducing predictive accuracy of the
$k{=}2$ horizon. Increasing $k$ or using Kalman filtering for
smoother velocity estimates is a direction for future work.

\subsection{Research Gap Resolution Summary}
Table~\ref{tab:gaps} summarises the improvement of each novel model
over the base paper's published CPF results.

\begin{table}[h]
\caption{Novel Model Gains vs.\ Base Paper CPF~\cite{xu2025}.
$\Delta CR$ = Capture Rate improvement in percentage points.}
\label{tab:gaps}
\centering
\renewcommand{\arraystretch}{1.2}
\begin{tabular}{l r r r r}
\toprule
\textbf{Model} &
\multicolumn{2}{c}{\textbf{Miami}} &
\multicolumn{2}{c}{\textbf{Test Maze}} \\
\cmidrule(lr){2-3} \cmidrule(lr){4-5}
& $\Delta CR$ & $\Delta\overline{CT}$ (s)
& $\Delta CR$ & $\Delta\overline{CT}$ (s) \\
\midrule
PI-PF    & $+14.2$~pp & $-2.89$ & $+23.0$~pp & $-2.40$ \\
RL-Q     & $+19.6$~pp & $-10.54$& $+47.6$~pp & $-22.38$ \\
APF-MCTS & $+19.8$~pp & $-11.04$& $+47.6$~pp & $-22.39$ \\
\bottomrule
\end{tabular}
\end{table}

%% ============================================================
\section{Conclusion}

We presented an independent, high-sample reproduction of the
state-of-the-art Composite Potential Field guard AI framework
(Xu~\& Verbrugge, AIIDE~2025) and verified all five baseline
models against published Table~3 metrics with $N{=}500$ Monte-Carlo
trials. We identified two fundamental scientific gaps in the base
paper---Reactive Pursuit Lag and Hand-Tuned Heuristic Weights---and
addressed them with three novel AI models.

Our results establish that:
\begin{enumerate}
  \item \textbf{PI-PF} eliminates reactive lag ($CR = 0.942$ on
        Miami, $+14.2$~pp over paper CPF, $B = 0.00$).
  \item \textbf{RL-Q} autonomously learns coordinated
        multi-agent flanking via Q-learning, achieving $CR = 0.996$
        on both topologies with $\overline{CT} \approx 21$~s---
        matching the theoretical Cheat upper bound and surpassing
        base paper CPF by $+47.6$~pp on the maze.
  \item \textbf{APF-MCTS} eliminates dead-end local minima via
        $H{=}3$ online tree search, achieving identical $CR = 0.996$
        with the lowest mean capture time and zero backtracking
        across all maps.
\end{enumerate}

Future work includes: (1) validating on NavMesh-based continuous
environments to close the Staleness fidelity gap; (2) extending
PI-PF to multi-guard coordinated interception using convex hull
trajectory prediction; and (3) replacing tabular Q-learning with
Deep Q-Networks (DQN) to support continuous state spaces and
partial observability.

%% ============================================================
\section*{Acknowledgments}
The authors acknowledge Kaijie Xu and Clark Verbrugge
(McGill University) for publishing the benchmark maps and
experimental protocol that enabled this reproducibility study.

%% ============================================================
\begin{thebibliography}{10}

\bibitem{xu2025}
K.~Xu and C.~Verbrugge,
``Generic Guard AI in Stealth Game with Composite Potential Fields,''
in \emph{Proc.\ 21st AAAI Int.\ Conf.\ Artif.\ Intell.\ Interactive
Digital Entertain.\ (AIIDE)}, 2025. arXiv:2508.18527.

\bibitem{alenezi2020}
A.~Al~Enezi and C.~Verbrugge,
``Dynamic Guard Patrol in Stealth Games,''
in \emph{Proc.\ 16th AAAI AIIDE}, 2020.

\bibitem{alenezi2021}
A.~Al~Enezi and C.~Verbrugge,
``Skeleton-Based Multi-Agent Opponent Search,''
in \emph{IEEE Conf.\ Games (CoG)}, 2021.

\bibitem{li2024}
X.~Li, J.~Wang, and Y.~Liu,
``Generative Subgoal Oriented Multi-Agent Reinforcement Learning
Through Potential Field,''
\emph{Neural Networks}, vol.~180, 2024.

\bibitem{isla2013}
D.~Isla,
``Third Eye Crime: Building a Stealth Game Around Occupancy Maps,''
in \emph{Proc.\ 9th AAAI AIIDE}, 2013.

\bibitem{hladky2008}
D.~Hladky and V.~Bulitko,
``An Evaluation of Opponent Modelling in One-on-One Games,''
in \emph{IEEE Symp.\ Comput.\ Intell.\ Games (CIG)}, 2008.

\bibitem{sutton2018}
R.~S.~Sutton and A.~G.~Barto,
\emph{Reinforcement Learning: An Introduction}, 2nd~ed.
Cambridge, MA: MIT Press, 2018.

\bibitem{khatib1986}
O.~Khatib,
``Real-Time Obstacle Avoidance for Mobile Robots and Manipulators,''
\emph{Int.\ J.\ Robot.\ Res.}, vol.~5, no.~1, pp.~90--98, 1986.

\bibitem{browne2012}
C.~B.~Browne \emph{et~al.},
``A Survey of Monte Carlo Tree Search Methods,''
\emph{IEEE Trans.\ Comput.\ Intell.\ AI Games}, vol.~4, no.~1,
pp.~1--43, 2012.

\bibitem{britton2022}
G.~Britton,
``Decision-Making Decisions: Artificial Intelligence in Computer
Games,''
\emph{Undergraduate Research Report}, 2022.

\bibitem{yannakakis2012}
G.~N.~Yannakakis,
``Game AI Revisited,''
in \emph{Proc.\ 9th Conf.\ Comput.\ Frontiers (CF)},
pp.~285--292, 2012.

\end{thebibliography}

\end{document}
